\section{Sequence of Play}\label{sec:4.0}

\ruleslabel{Cases}

\subsection{The Game-Turn}\label{sec:4.1}

Each game is played in seventeen Game-Turns, each of which consists of two
Player-Turns (one Yugoslav and the other Axis) plus three independent Stages.
The Player whose Player-Turn is in progress is termed the Phasing Player. On
Game-Turn 1, the game begins with the Yugoslav Player-Turn (the Special Events
Stage is skipped). On Game-Turn 17, the game ends at the conclusion of the
Victory Point Stage (the Axis Player-Turn and Terminal Stage are skipped).

\subsection{Game-Turn Sequence Outline}\label{sec:4.2}

Each Game-Turn must proceed strictly in the order outlined below.

\textbf{A. Special Events Stage}

\begin{enumerate}
  \item \textbf{Allied Progress Phase} (Game-Turn 6 and after): The Yugoslav
  Player rolls one die and applies the result to the Allied Progress Track (see
  10.1).

  \item \textbf{Weather Phase} (Game-Turns 2, 6, 10, and 14): The Yugoslav
  Player rolls a single die to determine whether drought is created in the
  current and three successive Game-Turns (see 11.2).

  \item \textbf{Tito Phase} (when eligible, through Game-Turn 14): The Axis
  Player rolls a single die to identify or locate Tito (see 9.41).

  \item \textbf{Axis Reinforcement Phase}: Axis reinforcements and replacements
  are placed on the map. Units are also upgraded and transfers may take place
  (see 13.0).

  \item \textbf{Chetnik Collaboration Phase} (Game-Turns 2 through 17): The
  Yugoslav Player rolls a single die for each Chetnik-occupied box, triangle,
  or circle to determine whether the Chetniks become pro-Yugoslav, pro-Axis,
  or neutral (see 7.2). During Game-Turns 9 through 13, the Yugoslav Player also
  determines whether Allied Support Withdrawal has occurred (see 7.26).

  \item \textbf{Italian Surrender Phase}: Italian units disband, defect, or
  remain in place during the Game-Turn of Italian Surrender. Operation
  Konstantin also takes place (see 10.3 and 10.36).

  \item \textbf{Axis Anti-Guerrilla Operations Phase} (Game-Turns 3 through
  14; see 8.4):
  \begin{description}[leftmargin=1.1em,style=nextline,font=\bfseries]
    \item[Planning Segment] The Axis Player declares zero, one, or two
    Anti-Guerrilla Operations. For each declared operation, the target Zone is
    secretly recorded and up to 7 divisions are chosen to participate.
    \item[Yugoslav Reaction Segment] The Yugoslav Player rolls one die; the
    result is the number of Yugoslav units that may conduct a bonus Movement
    Phase.
    \item[Deployment Segment] Participating Axis units are placed in the
    Mountains and/or Hide-aways of the target Zone.
    \item[Combat Segment] Axis units attack Yugoslav units in corresponding
    triangles or circles of the target Zone Display.
  \end{description}
\end{enumerate}

\textbf{B. Yugoslav Player-Turn}

\begin{enumerate}
  \item \textbf{Movement Phase}: The Yugoslav Player may move all eligible
  units (see 6.4) and receives predetermined reinforcements.
  \item \textbf{Combat Phase}: The Yugoslav Player attacks Axis units occupying
  corresponding locations in Objective or Zone Displays (see 8.0).
\end{enumerate}

\textbf{C. Victory Point Stage}

The Yugoslav Player accumulates or deducts Victory Points according to Section
14.0.

\textbf{D. Axis Player-Turn}

\begin{enumerate}
  \item \textbf{Movement Phase}: The Axis Player may move all eligible units
  (see 6.3).
  \item \textbf{Combat Phase}: The Axis Player attacks Yugoslav units occupying
  corresponding locations in Objective or Zone Displays (see 8.0).
\end{enumerate}

\textbf{E. Terminal Stage}

\begin{enumerate}
  \item \textbf{Guerrilla Reinforcement Phase}:
  \begin{description}[leftmargin=1.1em,style=nextline,font=\bfseries]
    \item[Recruitment Segment] Deploy guerrilla reinforcements created by the
    occupation of Objective or Zone Displays (see 7.43).
    \item[Tito Segment] Deploy Partisan reinforcements created by Tito's
    presence in a Zone (see 7.46).
    \item[Uprising Segment] From Game-Turn 3 onward, deploy guerrilla
    reinforcements created by successful Guerrilla Uprisings (see 7.47).
  \end{description}
  \item \textbf{Guerrilla Status Phase}: Determine whether Partisan and/or
  Chetnik units have achieved brigade or division strength (see 7.6).
  \item \textbf{Axis Anti-Guerrilla Operations Redeployment Phase}: Return Axis
  units that participated in Anti-Guerrilla Operations to the map (see 8.47).
  \item \textbf{Game-Turn Indication Phase}: Advance the Game-Turn marker one
  space on the Turn Record Track.
\end{enumerate}
